\documentclass[10pt,a4paper]{amsart}
\usepackage[margin=19mm]{geometry}
\usepackage{amsmath,amssymb,mathtools}
\usepackage[scr=rsfs,cal=euler]{mathalfa}
\usepackage{xfrac}
\usepackage[unicode,hidelinks,bookmarks=false]{hyperref}
\newcommand{\ie}{i.e.\ }
\newcommand{\cf}{cf.\ }
\newcommand{\resp}{respectively\ }
\newcommand{\Subsection}{\subsection}
\providecommand{\nobreakdash}{\nobreak\hbox{-}}
\setlength{\emergencystretch}{2em}
\multlinegap0pt
\usepackage{amssymb}
\usepackage{enumerate}
\newtheorem{lemma}{Lemma}
\newtheorem{theorem}{Theorem}
\title{On the Right Eigenvalues of the Quaternionic Matrix Polynomials}
\author{Ovaisa Jan, Idrees Qasim}
\date{Source: arXiv:2604.14813v1 (CC BY 4.0)}
\begin{document}
\maketitle

\noindent\emph{Source and licence.} On the Right Eigenvalues of the Quaternionic Matrix Polynomials. Version arXiv:2604.14813v1 (CC BY 4.0), distributed by arXiv under the Creative Commons Attribution 4.0 International licence, http://creativecommons.org/licenses/by/4.0/, as recorded in the arXiv licence metadata for this exact version and shown on the abstract page https://arxiv.org/abs/2604.14813v1. Full licence notice: \texttt{licenses/arxiv-2604.14813-CC-BY-4.0.txt}.\par\smallskip

\noindent\textit{Editorial scope.} Counted unit: the article's theorem 3.6 (source0.tex lines 490--505). The article's complete original proof of it is delivered with it (source0.tex lines 506--583, one \texttt{proof} environment of 78 lines). No mathematics is written, repaired or re-derived: the statement and the proof are the article's own lines, in the article's own notation, and no retained line is substituted or re-flowed.  Retained uncounted as the source-local context the proof needs: lines 265--285.  Nothing was omitted from any retained span, so the package carries no omission record.  No retained line contains a citation, so the wrapper carries no bibliography and no bibliography entry.  No retained line contains a figure environment, an image-inclusion command or drawing code, so no image is delivered and the package declares no asset dependency.  Editorial formatting only, in addition: an AMS \texttt{amsart} preamble that restates the article's own declarations and macros used by the retained lines, namely the environments \texttt{lemma}, \texttt{theorem} on separate counters, together with the standard packages those lines need.  The retained environments print in this document as Lemma 1, Theorem 1 in order of appearance, whereas the article prints the same items with the numbering of its own full document.  The complete native source of the article as distributed in the arXiv e-print for this exact version is bundled unchanged as \texttt{sources/198557/source0.tex}.
    \begin{lemma}{\label{lemma 3.2}}
    Let \( A = [a_{ij}] \in M_n(\mathbb{H}) \) be partitioned as
    \[
    A = \begin{bmatrix}
        A_{11} & A_{12} \\
        A_{21} & A_{22}
    \end{bmatrix},
    \]
    where \( A_{ij} \in M_{n_i \times n_j}(\mathbb{H}) \) is the \((i,j)\) block of \( A \) such that \( n_1 + n_2 = n \), where \( i,j \in \{1,2\} \). If
    \[
    \tilde{A} = \begin{bmatrix}
        \|A_{11}\|_2 & \|A_{12}\|_2 \\
        \|A_{21}\|_2 & \|A_{22}\|_2
    \end{bmatrix},
    \]
    then the following inequalities hold:
    \begin{itemize}
        \item \( \rho_r(A) \leq \rho_r(\tilde{A}) \).
        \item \( \|A\|_2 \leq \|\tilde{A}\|_2 \).
    \end{itemize}
\end{lemma}

\begin{theorem}{\label{theorem 3.6}}
     For every right eigenvalue $\xi$ of the quaternionic monic matrix polynomial $L(\xi)$ the following inequality holds:\\
    \begin{equation*}
|\xi| \le \left( 1+\frac{1}{2} \left[ \beta^{2}_{1} + \beta^{2}_{2} +\sqrt{(\beta^{2}_{1}+\beta^{2}_{2})^{2}+4\left(2\beta_{1}\beta_{2} + 1 \right) }\right] \right)^{\frac{1}{4}},
\end{equation*}\\
where \begin{multline*}
 \beta_{1}=\left( \frac{1}{2} \left[ \sum_{i=0}^{k-3}\left(\|A_{i}\|^{2}_{2}+\|B_{i}\|_{2}^{2}\right)+ \right. \right. \\
\left. \left.\sqrt{ \left( \sum_{i=0}^{k-3}\left(\|A_{i}\|^{2}_{2}-\|B_{i}\|_{2}^{2} \right)     \right)^{2}            +4\left|\left|\sum_{i=0}^{k-3}A_{i}{B_{i}^{H}}\right|\right|_{2}^{2} }~~\right]  \right)^{\frac{1}{2}},
\end{multline*}
\begin{multline*} 
\beta_{2}=\left( \frac{1}{2} \left[ \sum_{i=k-2}^{k-1}\left(\|A_{i}\|^{2}_{2}+\|B_{i}\|_{2}^{2}\right)+ \right. \right. \\
\left. \left. \sqrt{ \left( \sum_{i=k-2}^{k-1}\left(\|A_{i}\|^{2}_{2}-\|B_{i}\|_{2}^{2} \right) \right)^{2} +4\left\|\sum_{i=k-2}^{k-1}A_{i}{B_{i}^{H}}\right\|_{2}^{2} }~~\right] \right)^{\frac{1}{2}}. 
\end{multline*}


\end{theorem}

\begin{proof}
   Let
\[
C^{2}_{B} = \begin{bmatrix} N_{11} & N_{12} \\ N_{21} & N_{22} \end{bmatrix}
\] be the block matrix of $C^{2}_{L}$
where,
\[
N_{11} = 
\begin{bmatrix}
0 & 0 & I_n & \cdots & 0 \\
\vdots & \vdots  & \vdots &  &\vdots\\
0 & 0 & 0 & \cdots & I_n \\
0 & 0 & 0 & \cdots & 0 \\
0 & 0 & 0 & \cdots & 0 \\
\end{bmatrix}, \quad N_{12} = 
\begin{bmatrix}
0 &0 \\
\vdots&\vdots \\
I_n & 0\\ 0 & I_{n}
\end{bmatrix},
\]

\[
N_{21} = 
\begin{bmatrix}
-A_0 & -A_1 & \cdots & -A_{k-3}\\
B_0 & B_1 & \cdots & B_{k-3}
\end{bmatrix}, \quad N_{22} = \begin{bmatrix}
    -A_{k-2}&-A_{k-1}\\B_{k-2}&B_{k-1}\end{bmatrix}
\]
Since, $\|N_{11}\|_{2}=\|N_{12}\|_{2}=1.$
To find $\|N_{21}\|_{2},$ we compute  \[
A=N_{21}N_{21}^{H} = 
\begin{bmatrix}
\sum_{i=0}^{k-3} A_{i}A_{i}^{H}& -\sum_{i=0}^{k-3}A_{i}A_{i}^{H}  \\
 -\sum_{i=0}^{k-3} B_{i}A_{i}^{H}& \sum_{i=0}^{k-3} B_{i}B_{i}^{H}
\end{bmatrix}. \]
Thus, \[\tilde{A}= 
\begin{bmatrix}
\|\sum_{i=0}^{k-3} A_{i}A_{i}^{H}\|_{2}& -\|\sum_{i=0}^{k-3}A_{i}A_{i}^{H}\|_{2}  \\
 -\|\sum_{i=0}^{k-3} B_{i}A_{i}^{H}\|_{2}& \|\sum_{i=0}^{k-3} B_{i}B_{i}^{H}\|_{2}
\end{bmatrix}. \] By applying Lemma \ref{lemma 3.2}, we obtain 

\begin{align*}
\rho_{r}(A) &\le \rho_{r}(\tilde{A}) 
= \frac{1}{2}\Biggl[ \biggl( \Bigl\|\sum_{i=0}^{k-3} A_i A_i^{H}\Bigr\|_2 + \Bigl\|\sum_{i=0}^{k-3} B_i B_i^{H}\Bigr\|_2 \biggr) \\
&~~~~~\qquad + \sqrt{ \biggl( \Bigl\|\sum_{i=0}^{k-3} A_i A_i^{H}\Bigr\|_2 - \Bigl\|\sum_{i=0}^{k-3} B_i B_i^{H}\Bigr\|_2 \biggr)^2 + 4\Bigl\|\sum_{i=0}^{k-2} A_i B_i^{H}\Bigr\|_2 \Bigl\|\sum_{i=0}^{k-2} B_i A_i^{H}\Bigr\|_2 } \Biggr] \\
&~~~~~~~~~~~~\le \frac{1}{2}\Biggl[ \biggl( \sum_{i=0}^{k-3} \|A_i A_i^{H}\|_2 + \sum_{i=0}^{k-3} \|B_i B_i^{H}\|_2 \biggr) \\
&~~~~~\qquad + \sqrt{ \biggl( \sum_{i=0}^{k-3} \|A_i A_i^{H}\|_2 - \sum_{i=0}^{k-3} \|B_i B_i^{H}\|_2 \biggr)^2 + 4\Bigl\|\sum_{i=0}^{k-2} A_i B_i^{H}\Bigr\|_2 \Bigl\|\sum_{i=0}^{k-2} B_i A_i^{H}\Bigr\|_2 } \Biggr] \\
&~~~~~~~~~~~= \frac{1}{2}\Biggl[ \biggl( \sum_{i=0}^{k-3} \|A_i\|_2^2 + \sum_{i=0}^{k-3} \|B_i\|_2^2 \biggr) \\
&~~~~~~\qquad + \sqrt{ \biggl( \sum_{i=0}^{k-3} \|A_i\|_2^2 - \sum_{i=0}^{k-3} \|B_i\|_2^2 \biggr)^2 + 4\Bigl\|\sum_{i=0}^{k-2} A_i B_i^{H}\Bigr\|_2 \Bigl\|\sum_{i=0}^{k-2} A_i B_i^{H}\Bigr\|_2 } \Biggr] \\
&~~~~~~~~~~~= \frac{1}{2}\Biggl[ \biggl( \sum_{i=0}^{k-3} \|A_i\|_2^2 + \sum_{i=0}^{k-3} \|B_i\|_2^2 \biggr) \\
&~~~~~~\qquad + \sqrt{ \biggl( \sum_{i=0}^{k-3} \|A_i\|_2^2 - \sum_{i=0}^{k-3} \|B_i\|_2^2 \biggr)^2 + 4\Bigl\|\sum_{i=0}^{k-2} A_i B_i^{H}\Bigr\|_2^2 } \Biggr].
\end{align*}
 From the above inequalities, we can easily conclude 
     $\|N_{21}\|_{2}\le \beta_{1},$ \\
where  \begin{multline*}
 \beta_{1}=\left( \frac{1}{2} \left[ \sum_{i=0}^{k-3}\left(\|A_{i}\|^{2}_{2}+\|B_{i}\|_{2}^{2}\right)+ \right. \right. \\
\left. \left.\sqrt{ \left( \sum_{i=0}^{k-3}\left(\|A_{i}\|^{2}_{2}-\|B_{i}\|_{2}^{2} \right)     \right)^{2}            +4\left|\left|\sum_{i=0}^{k-3}A_{i}{B_{i}^{H}}\right|\right|_{2}^{2} }~~\right]  \right)^{\frac{1}{2}}.
\end{multline*}
Similarly, we obtain $\|N_{22}\|_{2}\le \beta_{2},$ where
\begin{multline*} 
\beta_{2}=\left( \frac{1}{2} \left[ \sum_{i=k-2}^{k-1}\left(\|A_{i}\|^{2}_{2}+\|B_{i}\|_{2}^{2}\right)+ \right. \right. \\
\left. \left. \sqrt{ \left( \sum_{i=k-2}^{k-1}\left(\|A_{i}\|^{2}_{2}-\|B_{i}\|_{2}^{2} \right) \right)^{2} +4\left\|\sum_{i=k-2}^{k-1}A_{i}{B_{i}^{H}}\right\|_{2}^{2} }~~\right] \right)^{\frac{1}{2}}. 
\end{multline*}
Now using Lemma {\ref{lemma 3.2}}, we have 

$$ \|C_{B}^{2}\|_2 \leq \|\tilde{C_{B}^{2}}\|_2. $$
This implies
\begin{align*}
   \|C_{B}^{2}\|^2_{2} &\le \left\| \begin{bmatrix} \|N_{11}\|_{2} &\|N_{12}\|_{2}\\ \|N_{21}\|_{2} & \|N_{22}\|_{2}\end{bmatrix} \right\|^2_{2} \\&=\left\| \begin{bmatrix} 1 &1\\ \|N_{21}\|_{2} & \|N_{22}\|_{2} \end{bmatrix} \right\|^2_{2}\\&\le \left\|\begin{bmatrix} 1 &1\\ \beta_{1} & \beta_{2} \end{bmatrix}\right\|^2_{2}.
\end{align*}
Thus, 
$$ \|C_{B}^{2}\|_{2} \le \left( 1+\frac{1}{2} \left[ \beta^{2}_{1} + \beta^{2}_{2} +\sqrt{(\beta^{2}_{1}+\beta^{2}_{2})^{2}+4\left(2\beta_{1}\beta_{2} + 1 \right) }\right] \right)^{\frac{1}{2}}.$$

\noindent Hence every right eigenvalues satisfies $$|\xi| \le \left( 1+\frac{1}{2} \left[ \beta^{2}_{1} + \beta^{2}_{2} +\sqrt{(\beta^{2}_{1}
+\beta^{2}_{2})^{2}+4\left(2\beta_{1}\beta_{2} + 1 \right) }\right] \right)^{\frac{1}{4}},$$
\end{proof}

\end{document}
